\section{Introduction}
\label{sec:intro}

\subsection{The bottleneck is not papers. It is environments.}

Life-science research produces an unprecedented volume of computationally
specified methods. Europe PMC offers programmatic access to over 37~million
publications, including 3.2~million fully open-access articles
\cite{europepmc}. PubMed Central, bioRxiv, and medRxiv provide additional APIs
for bulk retrieval. The raw material for constructing agent environments is
abundant and openly accessible.

The bottleneck lies in the \emph{transformation from paper to executable
environment}. An executable environment requires not only a code implementation
of the method but also a reliable evaluator that distinguishes genuine
methodological contributions from shortcuts. A benchmark study of biomedical
paper reproduction (ClaroAI-Bench) found that \emph{environment reconstruction
is the dimension with the highest disagreement among evaluators}, exceeding
even code availability as a source of variance \cite{claroaibench}. Papers with
accessible data and code achieve reproducibility scores $2.9\times$ higher than
those with restricted access, but a substantial fraction of published work does
not meet this threshold.

\subsection{\harbor\ provides the execution substrate.}

The \harbor\ framework, developed by Laude Institute (creators of
Terminal-Bench), provides a standardized substrate for agent evaluation in
containerized environments \cite{harbor}. \harbor\ is the official harness for
Terminal-Bench~2.0 \cite{terminalbench} and supports arbitrary agents including
Claude Code and open-source models. In the life sciences, \harbor-compatible
task sets already exist:

\begin{itemize}
  \item \textbf{DrugDiscoveryBench} (Scale AI): 82 \harbor-formatted tasks
        spanning target identification, patent mining, cheminformatics, and
        structural reasoning, each graded by an LLM judge against
        expert-authored rubrics \cite{drugdiscoverybench}.
  \item \textbf{Terminal-Bench-Science}: 70 expert-curated research workflow
        tasks covering life sciences, physics, and earth sciences
        \cite{tbscience}.
  \item \textbf{BioAgent Gym}: a biomedical agent repository with a native
        \harbor\ leaderboard \cite{bioagentgym}.
  \item \textbf{RSI-Exam}: evaluation infrastructure that runs directly on
        \harbor\ for assessing long-horizon self-improvement and generalization
        to unseen data \cite{rsiexam}.
\end{itemize}

These task sets are \emph{manually curated by domain experts}. The missing
component is automation: a pipeline that can generate \harbor\ tasks directly
from the published literature.

\subsection{Contribution}

We propose \sys, comprising:

\begin{enumerate}
  \item A \textbf{multi-agent pipeline} that ingests open-access life-science
        papers and outputs \harbor-formatted executable environments with
        automated evaluators.
  \item A \textbf{four-stage validation protocol} that progresses from
        environment quality verification through weight-free self-improvement
        and harness optimization to post-training with held-out benchmarks.
  \item A \textbf{scoped rollout plan} focused on computational biology and
        bioinformatics, the subfields where papers are most readily convertible
        to executable tasks.
\end{enumerate}

A fourth contribution is pragmatic but load-bearing: we release the pipeline as
a dependency-free open-source framework in which every environment must pass
four verification gates before it can be packaged (\S\ref{sec:artifact}). The
framework's structural claims are measured by its own test suite
(Appendix~\ref{app:claims}); its empirical claims are explicitly registered as
not-yet-measured.
